\scanpage{116}
\[
\deleted{(\alpha,\beta)\longmapsto}
\]
\[
\begin{aligned}
F_M&:C_M\times C_M\unsure{-|\longrightarrow}D_M\\
F_M(\alpha,\beta)&=\operatorname{tr}(\alpha\beta).1-[\alpha,\beta]
\end{aligned}
\]
Compatible avec l'hom
\[
G_m\times G_m\xrightarrow{(\lambda,\mu)\unsure{-|\mapsto}\lambda\mu}G_m
\]
sur les groupes d'homothéties. \deleted{\unsure{Passant aux} cônes épointés, cela se traduit par un hom de $G_m$-torseurs}
\[
\deleted{C_M^{\unsure{\times|\wedge}}\times C_M^{\unsure{\times|\wedge}}/\underset{\text{\unsure{agissant|opérant|épointé} antidiagonalement}}{G_m}\ \underset{\unsure{-}}{\sim}\ D_M^{\unsure{\times|\wedge}}}
\]

Le fait que l'on s'envoie bien dans $D_M$ provient de la formule
\[
\begin{aligned}
&\alpha,\beta\ \text{\unsure{str.|étr.|étant} nilpotents}\qquad(\text{i.e. }\det\alpha=\det\beta=\operatorname{tr}\alpha=\operatorname{tr}\beta=0)\\
&\Longrightarrow q_M\bigl([\alpha,\beta]/\varphi(\alpha,\beta)\bigr)=1\\
&\hspace{2em}\text{i.e.}\quad\det([\alpha,\beta])+\unsure{\operatorname{tr}|\varphi}(\alpha\beta)^2=0
\end{aligned}
\]
i.e.
\[
q_M[\alpha,\beta]=\varphi(\alpha,\beta{\unsure{]|)}}^2
\]
ce qui provient en effet de la formule générale (valable pour $\alpha,\beta\in\Gamma\unsure{\gamma|\mathfrak g}M$ sans plus)
\[
q_M\bigl([\alpha,\beta]\bigr)=-4q(\alpha)q(\beta)+\varphi(\alpha,\beta)^2
\]
% archive sheet number: 115
